% GENERATED FILE. Edit canonical lessons, then run npm run generate:books.
% canonical-source-hash: fnv1a64:9bb89c4fff54d451
% canonical-lessons: ZH-C209-guiding, ZH-C209-anquan, ZH-C209-baochi, ZH-C209-darao, ZH-C209-yunxu

\chapter{Rules and Notices}
\label{ch:zh-rules-and-notices}
\hlchaptermodality{209}

\begin{chapteropening}
\textbf{By the end of this chapter:} \emph{I can read the rules on a notice: what is safe, keeping quiet, not disturbing, and what is not allowed.}

Complete the last lesson of chapter 209: I can read the rules on a notice: what is safe, keeping quiet, not disturbing, and what is not allowed.
\end{chapteropening}

\section[guīdìng]{guīdìng — rule; regulation}

\label{lesson:ZH-C209-guiding}



\begin{quote}
Before the new one: say the Mandarin for tool, then the Mandarin for battery.
\end{quote}

\subsection*{You'll want to know: guīdìng}
\emph{guīdìng} — \textquotedblleft{}rule; regulation\textquotedblright{}.

\begin{quote}
\emph{zhè shì xuéxiào de guīdìng} — That's the school's rule.
\end{quote}

Its characters come later; for now it is a word for the ear.

\begin{cousinweb}
\emph{guīdìng} is built from \emph{guī} 'a pair of compasses, a rule' and \emph{dìng} 'to settle', both inherited from Middle Chinese, \emph{kjwie} and \emph{dengH} (Guangyun, in Baxter's transcription), which glosses the first 'round' and the second 'calm'.
\end{cousinweb}

\begin{grammarlens}[title={What you've built}]
One more word for reading notices and messages.
\end{grammarlens}

\subsection*{Your turn}
\begin{itemize}
\raggedright
  \item \emph{Say it:} \emph{guīdìng}
  \item \emph{Say it:} \emph{guīdìng}, once more
  \item \emph{Say it:} \emph{zhè shì xuéxiào de guīdìng} — That's the school's rule.
  \item \emph{Recall:} say the Mandarin for tool, then the Mandarin for battery, then say \emph{guīdìng} again
\end{itemize}

\begin{tcolorbox}[breakable,colback=teal!4,colframe=teal!35!black,title={Before you move on}]
What does guīdìng mean? (\textquotedblleft{}Rule; regulation\textquotedblright{}.) Say it once more.
\end{tcolorbox}
\section[ānquán]{ānquán — safe; safety}

\label{lesson:ZH-C209-anquan}



\begin{quote}
Before the new one: say the Mandarin for battery, then the Mandarin for regulation.
\end{quote}

\subsection*{You'll want to know: ānquán}
\emph{ānquán} — \textquotedblleft{}safe; safety\textquotedblright{}.

\begin{quote}
\emph{wǎnshang yí ge rén zǒulù bù ānquán} — It isn't safe to walk on your own at night.
\end{quote}

Its characters come later; for now it is a word for the ear.

\begin{cousinweb}
\emph{ānquán} is built from \emph{ān} 'calm' and \emph{quán} 'whole', both inherited from Middle Chinese, \emph{'an} and \emph{dzjwen} (Guangyun, in Baxter's transcription).
\end{cousinweb}

\begin{grammarlens}[title={What you've built}]
One more word for reading notices and messages.
\end{grammarlens}

\subsection*{Your turn}
\begin{itemize}
\raggedright
  \item \emph{Say it:} \emph{ānquán}
  \item \emph{Say it:} \emph{ānquán}, once more
  \item \emph{Say it:} \emph{wǎnshang yí ge rén zǒulù bù ānquán} — It isn't safe to walk on your own at night.
  \item \emph{Recall:} say the Mandarin for battery, then the Mandarin for regulation, then say \emph{ānquán} again
\end{itemize}

\begin{tcolorbox}[breakable,colback=teal!4,colframe=teal!35!black,title={Before you move on}]
What does ānquán mean? (\textquotedblleft{}Safe; safety\textquotedblright{}.) Say it once more.
\end{tcolorbox}
\section[bǎochí]{bǎochí — to keep up; to keep as it is}

\label{lesson:ZH-C209-baochi}



\begin{quote}
Before the new one: say the Mandarin for regulation, then the Mandarin for safe.
\end{quote}

\subsection*{You'll want to know: bǎochí}
\emph{bǎochí} — \textquotedblleft{}to keep up; to keep as it is\textquotedblright{}.

\begin{quote}
\emph{qǐng bǎochí ānjìng} — Please keep quiet.
\end{quote}

Its characters come later; for now it is a word for the ear.

\begin{cousinweb}
\emph{bǎochí} is built from \emph{bǎo} 'to guard' and \emph{chí} 'to hold', both inherited from Middle Chinese, \emph{pawX} and \emph{dri} (Guangyun, in Baxter's transcription), which glosses the first 'to keep safe' and the second 'to grasp'.
\end{cousinweb}

\begin{grammarlens}[title={What you've built}]
One more word for reading notices and messages.
\end{grammarlens}

\subsection*{Your turn}
\begin{itemize}
\raggedright
  \item \emph{Say it:} \emph{bǎochí}
  \item \emph{Say it:} \emph{bǎochí}, once more
  \item \emph{Say it:} \emph{qǐng bǎochí ānjìng} — Please keep quiet.
  \item \emph{Recall:} say the Mandarin for regulation, then the Mandarin for safe, then say \emph{bǎochí} again
\end{itemize}

\begin{tcolorbox}[breakable,colback=teal!4,colframe=teal!35!black,title={Before you move on}]
What does bǎochí mean? (\textquotedblleft{}To keep up; to keep as it is\textquotedblright{}.) Say it once more.
\end{tcolorbox}
\section[dǎrǎo]{dǎrǎo — to disturb}

\label{lesson:ZH-C209-darao}



\begin{quote}
Before the new one: say the Mandarin for safe, then the Mandarin for to keep up.
\end{quote}

\subsection*{You'll want to know: dǎrǎo}
\emph{dǎrǎo} — \textquotedblleft{}to disturb\textquotedblright{}.

\begin{quote}
\emph{qǐng bú yào dǎrǎo wǒ} — Please don't disturb me.
\end{quote}

Its characters come later; for now it is a word for the ear.

Two third tones meet in this word: dǎ is said with a rising tone before rǎo, and the spelling keeps both marks.

\begin{cousinweb}
\emph{dǎrǎo} is built from \zh{打} 'to strike' and \emph{rǎo} 'to trouble', both inherited from Middle Chinese, \emph{taengX} and \emph{nyewX} (Guangyun, in Baxter's transcription), which glosses the second 'to throw into disorder' and 'to vex'.
\end{cousinweb}

\begin{grammarlens}[title={What you've built}]
One more word for reading notices and messages.
\end{grammarlens}

\subsection*{Your turn}
\begin{itemize}
\raggedright
  \item \emph{Say it:} \emph{dǎrǎo}
  \item \emph{Say it:} \emph{dǎrǎo}, once more
  \item \emph{Say it:} \emph{qǐng bú yào dǎrǎo wǒ} — Please don't disturb me.
  \item \emph{Recall:} say the Mandarin for safe, then the Mandarin for to keep up, then say \emph{dǎrǎo} again
\end{itemize}

\begin{tcolorbox}[breakable,colback=teal!4,colframe=teal!35!black,title={Before you move on}]
What does dǎrǎo mean? (\textquotedblleft{}To disturb\textquotedblright{}.) Say it once more.
\end{tcolorbox}
\section[yǔnxǔ]{yǔnxǔ — to allow; to permit}

\label{lesson:ZH-C209-yunxu}



\begin{quote}
Before the new one: say the Mandarin for to keep up, then the Mandarin for to disturb.
\end{quote}

\subsection*{You'll want to know: yǔnxǔ}
\emph{yǔnxǔ} — \textquotedblleft{}to allow; to permit\textquotedblright{}.

\begin{quote}
\emph{zhèr bù yǔnxǔ xīyān} — Smoking isn't allowed here.
\end{quote}

Its characters come later; for now it is a word for the ear.

Two third tones meet in this word: yǔn is said with a rising tone before xǔ, and the spelling keeps both marks.

\begin{cousinweb}
\emph{yǔnxǔ} is built from \emph{yǔn} and \emph{xǔ}, both inherited from Middle Chinese, \emph{ywinX} and \emph{xjoX} (Guangyun, in Baxter's transcription), which glosses the first 'true' and the second 'to permit'.
\end{cousinweb}

\begin{grammarlens}[title={What you've built}]
One more word for reading notices and messages.
\end{grammarlens}

\subsection*{Your turn}
\begin{itemize}
\raggedright
  \item \emph{Say it:} \emph{yǔnxǔ}
  \item \emph{Say it:} \emph{yǔnxǔ}, once more
  \item \emph{Say it:} \emph{zhèr bù yǔnxǔ xīyān} — Smoking isn't allowed here.
  \item \emph{Recall:} say the Mandarin for to keep up, then the Mandarin for to disturb, then say \emph{yǔnxǔ} again
\end{itemize}

\begin{tcolorbox}[breakable,colback=teal!4,colframe=teal!35!black,title={Before you move on}]
What does yǔnxǔ mean? (\textquotedblleft{}To allow; to permit\textquotedblright{}.) Say it once more.
\end{tcolorbox}
